\documentclass[12pt]{article}

\usepackage{graphicx}
\usepackage{url}
\usepackage{hyperref}
\usepackage{verbatim}
\usepackage[title,titletoc,toc]{appendix}
\usepackage{wasysym}

\renewcommand{\topfraction}{1.0}
\renewcommand{\bottomfraction}{1.0}
\renewcommand{\textfraction}{0.0}
\renewcommand{\floatpagefraction}{1.0}
\renewcommand{\dbltopfraction}{1.0}


\begin{document}

\begin{titlepage}

\begin{center}
{\large STAT 6910 Advanced R} \\[4cm]

{\LARGE \bf Reproducible Research (RR)} \\[1cm]
by \\[0.5cm]
{\bf J\"urgen Symanzik} \\[2.5cm]
{\bf Date:} \today \\[2cm]

UTAH STATE UNIVERSITY \\[0.5cm]
Logan, UT \\[0.5cm]
Fall 2013 \\[0.5cm]
\end{center}

\thispagestyle{empty}
\vfill
\end{titlepage}


\newpage 


\pagenumbering{roman}

\tableofcontents


\newpage

\pagenumbering{arabic}


\section{Basics about \LaTeX\ via RStudio}

First, notice that you can open a \LaTeX\ .tex file from within RStudio.
This file is recognized as a .tex file --- and you can translate it 
via the Compile PDF option from within RStudio!

When working with .tex in RStudio, you have options to introduce
some \LaTeX\ formatting via highlighting and mouse clicke. Let's
try this for the words bold, italic, and typewriter.


Now, we want a new subsection

Notice what is being produced here!


Now, we want a new subsection

Do this again --- but remove the $*$ after \verb|\subsection|.


And what about a bullet list hereafter?


\subsection{Math Symbols}

Let's try a few math symbols: $\sigma, \lambda,
s_1 = \sum_{i=1}^n i^2, \infty, \sqrt{x}, \bar{x}$ and a \smiley.


Now, do the sum in a math environment:
\[
s_1 = \sum_{i=1}^n i^2
\]


\section{\LaTeX\ Documentation}

Numerous web sites provide information about \LaTeX,
including tutorials and help pages. Here are some:

\begin{itemize}
\item \url{http://www.latex-project.org/}
\item \url{http://en.wikibooks.org/wiki/LaTeX}
\item \url{http://miktex.org/}
\item \url{http://www.artofproblemsolving.com/Wiki/index.php/LaTeX:Symbols}
\end{itemize}

If you need help on a specific latex topic, google for
\verb|latex| and the keyword or command you are interested in.



\newpage


\section{Using Sweave and knitr in RStudio}

Detailed information can be found at
\url{http://www.rstudio.com/ide/docs/authoring/overview}.

Basically, you need to start at \textbf{File $->$ New $->$ R Sweave}.
Let's do so.


If you want to use an existing \verb|.tex| file, you must change the file
extension to \verb|.Rnw| and add the line
\begin{verbatim}
\SweaveOpts{concordance=TRUE}
\end{verbatim}
after the 
\begin{verbatim}
\begin{document}
\end{verbatim}
line.
 

\end{document}

